\documentclass[11pt]{article}
\usepackage{amsmath,amssymb,amsthm,geometry,hyperref}
\geometry{margin=25mm}
\newtheorem{lemma}{Lemma}
\newtheorem{proposition}[lemma]{Proposition}
\newcommand{\N}{\mathrm N}
\newcommand{\OO}{\mathcal O}
\newcommand{\rad}{\operatorname{rad}}
\title{A reduction of the joint inverse and plain witness moment}
\author{}
\date{8 October 2026}
\begin{document}
\maketitle

\paragraph{Status.}
The coefficient identities, resonant-pair count, complete-residue mean,
and diagonal-strip reductions below are proved.
The desired bound for the remaining physical row sum is open.
No new zero-free boundary follows from this note.
The pinned upstream revision is
\begin{center}
\texttt{adc7f1241b42e322a6451854ab7e4b4c146bf78a}.
\end{center}
The Lean file \path{lean/JointMomentReduction.lean} checks the complementary
coefficient identity, finite pair-kernel expansion, and diagonal-strip
inequality. The ideal-counting and analytic arguments are not formalized.

\paragraph{The actual product.}
Put $D=U^r$, $N=U^m$, $D_*=U^t$, and $X=DN=U^s$.
Let $V$ be the source cutoff, equal to one on $[0,1]$ and zero on
$[2,\infty)$. Let $W$ be the positive annular profile.
The argument also permits two fixed smooth profiles $W_I,W_P$ supported
in $[a,b]$, with $0<a<b$ fixed. The source character is
\[
 \psi_u(K)=\eta(K)\,
   \operatorname{idealRowHom}(m_0^6 f_0^4 u,K).
\]
The symbols $m_0,f_0$ denote the fixed row data rather than the length $m$.
Source prime deletions are part of this common character.
For the same real shifts $\sigma,\omega$ in both factors,
\begin{align*}
 I_u&=D^{-1/2}\sum_d\mu(d)\psi_u(d)V(\N d/D_*)
      W_I(\N d/D)(\N d/D)^{-\sigma+i\omega},\\
 P_u&=N^{-1/2}\sum_n\psi_u(n)
      W_P(\N n/N)(\N n/N)^{-\sigma+i\omega}.
\end{align*}
All ideals in this note are nonzero integral ideals of the Eisenstein ring.
The sums displayed here are finite. Multiplicativity, including zeros, gives
\begin{equation}\label{eq:product}
 I_uP_u=X^{-1/2}\sum_K c_K\psi_u(K)(\N K/X)^{-\sigma+i\omega},
 \qquad
 c_K=\sum_{dn=K}\mu(d)V(\N d/D_*)W_I(\N d/D)W_P(\N n/N).
\end{equation}
Its support is $a^2X\le\N K\le b^2X$.
The shifts are common within a row. Actual detector shifts
$\sigma_u,\omega_u$ need not be common across rows.

On a compact short crossing strip with $r\le t-\kappa$, $\kappa>0$,
the inverse cutoff is exactly one once $bU^{-\kappa}\le1$.
Indeed $\N d\le bU^r\le U^t$ throughout its support.
The target there is therefore a pure dyadic M\"obius and plain mixed
moment. Its coefficient does not receive extra cancellation from the
truncation boundary.
For four pairwise distinct prime ideals $d_1,d_2,n_1,n_2$ in the
respective annuli, the inverse signs satisfy
$\mu(d_1)\mu(d_2)=+1$ and the radical modulus has norm comparable to
$(DN)^2$. This sector prevents an argument that assumes either a shortened
modulus or oscillating inverse signs in every remaining configuration.

\paragraph{Common derivative marks.}
For $j+k=\ell$, differentiating the product in its common shifts multiplies
the summand in \eqref{eq:product} by
$(-1)^j i^k\log(\N K/X)^\ell$.
Thus all marks of total order at most two are bounded on the product annulus.
Independent differentiation of the two factors instead inserts separate
logarithms and does not preserve the complete divisor cancellation below.
At fixed $D,N,D_*$, a two-dimensional Sobolev argument transports uniform
fixed-shift energies for $F,\partial_\sigma F,\partial_\omega F,
\partial_\sigma\partial_\omega F$ to arbitrary rowwise shifts
in $[0,1]\times[-H,H]$, with a polynomial loss in $1+H$.
No fixed-shift kernel identity is applied directly to a moving-shift sum.

\paragraph{Complementary coefficients and their limitation.}
For $K\ne1$, the complete divisor identity gives the exact equality
\begin{equation}\label{eq:complement}
 A_{D_*}(K):=\sum_{dn=K}\mu(d)V(\N d/D_*)
 =-\sum_{dn=K}\mu(d)(1-V(\N d/D_*)).
\end{equation}
For an aggregate annular detector
\[
 T_u(X)=X^{-1/2}\sum_K\psi_u(K)A_{D_*}(K)
       W(\N K/X)(\N K/X)^{-\sigma_u+i\omega_u},
\]
whose support excludes the unit, \eqref{eq:complement} yields
\begin{equation}\label{eq:aggregate}
 T_u(X)=-\sum_{\N n\le bX/D_*}\frac{\psi_u(n)}{\sqrt{\N n}}
        I_u(X/\N n\,\,W_n),\qquad
 W_n(y)=W(y)\bigl(1-V(Xy/(\N nD_*))\bigr).
\end{equation}
The inverse polynomial on the right has scale $X/\N n$ and the same
rowwise shifts. Smooth seminorms of $W_n$ are uniform in $n,X,D_*$.
Derivatives of $V$ have fixed compact transition support.
Common marks in \eqref{eq:aggregate} remain powers of $\log y$.

The existing inverse amplification energy, for $s\ge t\ge1$, gives
\[
 \|I(X/\N n)\|_{\ell^2(u)}
 \ll (1+H)^{J/2}U^{1/12}(X/\N n)^{5/12}U^\rho.
\]
Minkowski and ideal counting then give
\begin{equation}\label{eq:aggregate-bound}
 \sum_u|T_u(X)|^2
 \ll (1+H)^J U^{1/6+s-t/6+\rho},
\end{equation}
because $\sum_{\N n\le Y}(\N n)^{-11/12}\ll Y^{1/12}$.
At $s=t$, this exponent equals $1+5(t-1)/6$.
The complementary variable has bounded norm there, so orthogonality only
in that variable cannot supply a power saving.
If a common prime-deletion mask kills every nonunit $n$ allowed in
\eqref{eq:aggregate}, only the unit survives. This conditional collapse
does not apply automatically to every source ray label.

An individual dyadic coefficient $c_K$ need not cancel or have shorter
support. For separated annuli, $K=pq$ with $\N p\asymp D$ and
$\N q\asymp N$ can have only $d=p$ in \eqref{eq:product}, giving
a nonzero coefficient of size comparable to one at norm $DN$.

\paragraph{The exact physical kernel.}
At a fixed pair of common shifts define
\[
 b_K=c_K(\N K/X)^{-\sigma+i\omega},\qquad
 \mathcal K_U(K,L)=\sum_{u\in\mathcal F_U}\psi_u(K)\overline{\psi_u(L)}.
\]
Here $\mathcal F_U$ is the actual sixth-power-free physical row family,
or a specified subset, with row norm at most $U$.
The energy is exactly
\begin{equation}\label{eq:kernel}
 \sum_{u\in\mathcal F_U}|I_uP_u|^2
 =X^{-1}\sum_{K,L}b_K\overline{b_L}\mathcal K_U(K,L).
\end{equation}
At a good prime $P$, its local row factor on units is the canonical sextic
character to exponent $v_P(K)-v_P(L)$ modulo six. Its value at zero is
zero whenever $P\mid KL$, including exponent zero modulo six.
Consequently the union zero mask is $\rad(KL)$, not just the conductor
of the reduced quotient character.

For $q=\N P$, let $\chi_P$ be the actual sextic character and
$\theta_P$ the chosen nontrivial additive character. Its finite Fourier
factor is
\[
 G_P(e,h)=\sum_{z\in\OO/P}\chi_P^e(z)\theta_P(hz).
\]
The zeroth power retains the zero extension.
Then $G_P(0,0)=q-1$, $G_P(e,0)=0$ for $e\not\equiv0\pmod6$,
and $G_P(0,h)=-1$ for $h\ne0$.
For $e\not\equiv0$ and $h\ne0$,
\[
 G_P(e,h)=\chi_P^e(h)^{-1}\tau_P(e),\qquad |\tau_P(e)|^2=q.
\]
CRT gives a product of these factors, with the actual trace-coordinate
scalings. On nonresonant pairs its global zero Fourier mode vanishes.
This does not imply that \eqref{eq:kernel} vanishes on a physical annulus.
For generic coprime squarefree $K,L$, the radical modulus has norm
$Q\asymp X^2$. Poisson completion has dual volume $Q/U$.
Taking absolute values costs approximately $\sqrt Q=X$, already larger
than $U$ when $s>1$.

\begin{lemma}[Resonant pairs]\label{lem:resonant-count}
Fix a finite set $\Sigma$ of prime ideals. For $\N K,\N L\le BX$,
suppose $v_P(K)\equiv v_P(L)\pmod6$ for every $P\notin\Sigma$.
The number of ordered pairs is at most
\[
 C_B X\zeta_{\OO}(3)^2
       \prod_{P\in\Sigma}(1-(\N P)^{-1/2})^{-2}.
\]
\end{lemma}
\begin{proof}
Write uniquely $K=K_\Sigma A B_1^6$ and $L=L_\Sigma A B_2^6$,
where $A$ is sixth-power-free outside $\Sigma$ and all three ideals
$A,B_1,B_2$ are supported outside $\Sigma$.
For fixed $K_\Sigma,L_\Sigma,B_1,B_2$, ideal counting bounds the number
of $A$ by a constant times
\[
 \frac{X}{\max(\N K_\Sigma(\N B_1)^6,\N L_\Sigma(\N B_2)^6)}
 \le\frac{X}{\sqrt{\N K_\Sigma\N L_\Sigma}(\N B_1)^3(\N B_2)^3}.
\]
Sum the latter bound. The $B_1,B_2$ sums converge to at most
$\zeta_{\OO}(3)^2$. Each fixed-supported ideal sum is
$\prod_{P\in\Sigma}(1-(\N P)^{-1/2})^{-1}$.
\end{proof}

If $\Sigma$ varies with arithmetic height, this Euler factor must be
retained. It is at most $(\N\rad\Sigma)^4$, since $\N P\ge2$.
Thus a stated polynomial bound for $\N\rad\Sigma$ supplies a polynomial
height loss. Primes dividing $m_0f_0$ can often be removed entirely,
since their zero masks annihilate ideals containing them.

\begin{proposition}[Absolute resonant energy]\label{prop:resonant-energy}
Assume $|\mathcal F_U|\ll U$, $s$ is bounded, and the profile sup norms
are bounded. For any $\rho>0$, the absolute contribution to
\eqref{eq:kernel} of the pairs in Lemma~\ref{lem:resonant-count} is
\[
 \ll (1+H)^J U^{1+\rho},
\]
with the explicit fixed-supported Euler factor if needed.
The same bound holds for common derivative marks of total order at most two
and for arbitrary rowwise shifts in the stated bounded real-shift interval.
\end{proposition}
\begin{proof}
The coefficient $|c_K|$ is at most a profile constant times the number
of squarefree ideal divisors of $K$, hence $\ll_\epsilon(\N K)^\epsilon$.
The real shifts and common logarithmic marks are bounded on the product
annulus. The norm of every row character is at most one.
Lemma~\ref{lem:resonant-count}, $|\mathcal F_U|\ll U$, and the
normalization $X^{-1}$ give $\ll U X^{2\epsilon}$.
Choose $\epsilon$ after the bounded range of $s$ and the reserved $\rho$.
\end{proof}

\begin{proposition}[Complete-residue special case]\label{prop:complete}
Average \eqref{eq:product} over every residue class modulo the product
of all varying good primes in its ideal support, with each residue class
given equal weight. Fix the common shifts across this average.
Then its normalized mean-square is $\ll X^\rho$, including common marks
of total order at most two and the fixed-supported factor above.
\end{proposition}
\begin{proof}
CRT and the vanishing local zero Fourier mode eliminate every nonresonant
pair. The norm of each surviving normalized complete pair mean is at most
one. Apply Lemma~\ref{lem:resonant-count} and the coefficient bound in
Proposition~\ref{prop:resonant-energy}, then divide by $X$.
\end{proof}
This is a proved restricted averaging result. It gives no discrepancy
estimate comparing the complete average with the sixth-power-free annulus,
especially when its modulus has norm much larger than $U$.

\paragraph{Two further off-diagonal strips are controlled.}
Write
\[
\begin{aligned}
 a_d&=\mu(d)V(\N d/D_*)W_I(\N d/D)(\N d/D)^{-\sigma+i\omega},\\
 b_n&=W_P(\N n/N)(\N n/N)^{-\sigma+i\omega}.
\end{aligned}
\]
Before collecting $K=dn$, the exact Type II sum is
\begin{equation}\label{eq:typeII}
 (DN)^{-1}\sum_{d_1,d_2,n_1,n_2}
 a_{d_1}\overline{a_{d_2}}b_{n_1}\overline{b_{n_2}}
 \mathcal K_U(d_1n_1,d_2n_2).
\end{equation}
The part with $d_1=d_2$ equals
\[
 D^{-1}\sum_{u,d}|a_d|^2|\psi_u(d)|^2|P_u|^2
 \ll\sum_u|P_u|^2.
\]
The part with $n_1=n_2$ similarly is at most a constant times
$\sum_u|I_u|^2$. These identities also permit moving common shifts,
because the coefficient square sums are uniformly $O(D)$ and $O(N)$.
Under the source marked marginal moment hypotheses both bounds are
$\ll(1+H)^J U^{1+\rho}$ in the crossing strip $r\le1$, $m\le1/2$.
For the inverse factor use its existing exponent
$\max(1,(1+5r)/6)=1$. For the plain factor a fourth-moment bound
$\sum_u|P_u|^4\ll U^{1+\rho}$ and $|\mathcal F_U|\ll U$ suffice
by Cauchy--Schwarz.

For a common mark $\ell\le2$, expand
$\log(\N(dn)/(DN))^\ell=
 (\log(\N d/D)+\log(\N n/N))^\ell$.
Cauchy--Schwarz over the at most three binomial terms reduces the same
strips to marked marginal energies. All factor logarithms are bounded
on their individual annuli.
The intersection $d_1=d_2,n_1=n_2$ is nonnegative. Inclusion and exclusion
therefore bound the total energy from above by these two controlled strips
plus the real part of the doubly off-diagonal sum.
This controls some nonresonant pairs without termwise kernel bounds.
It is conditional only on the already stated marginal source moments,
not on a new joint estimate.

\paragraph{What a pairwise square-root estimate would still miss.}
Even suppose, more strongly than the available completion argument,
that every remaining incomplete pair kernel satisfied
$|\mathcal K_U(K,L)|\ll U^{1/2+\rho}$.
There are $O(X^{2+\rho})$ weighted pairs. Absolute summation in
\eqref{eq:typeII} would then give only
$\ll U^{s+1/2+\rho}$ after normalization by $X$.
The same estimate follows from Cauchy--Schwarz using only a mean-square
pair-kernel bound $\sum_{K,L}|\mathcal K_U(K,L)|^2\ll X^2U^{1+\rho}$
and the coefficient square sum $\sum_K|c_K|^2\ll X^{1+\rho}$.
To reach the desired exponent requires additional signed correlation
cancellation of $U^{s/6+1/3+\eta}$ relative to those bounds.
At the displayed crossing this missing saving is about
$U^{0.520564+\eta}$.
A finite-field four-shift square-root bound alone supplies no such
global coefficient cancellation or transfer across varying moduli.

\paragraph{An exact Cauchy reorganization.}
At fixed common shifts put
\[
 A_u=\sum_d a_d\psi_u(d),\qquad B_u=\sum_n b_n\psi_u(n),
 \qquad J_u=\sum_d|a_d\psi_u(d)|^2,
 \quad L_u=\sum_n|b_n\psi_u(n)|^2.
\]
Before removing resonance, the doubly off-diagonal sum has the exact form
\begin{equation}\label{eq:covariance}
 \mathcal E_{\mathrm{both\ off}}
   =(DN)^{-1}\sum_u (|A_u|^2-J_u)(|B_u|^2-L_u).
\end{equation}
Expand both parentheses to verify this equality.
The four terms are the full energy, the two strips, and their intersection.
This subtraction is checked in Lean as
\path{express_double_strip_subtraction_as_centered_row_product}.
For common derivative marks, its polarized version is a finite sum of
centered cross-products between the marked factors, obtained by the same
binomial expansion. It does not introduce independent row shifts.

Cauchy over the short coefficient pair gives another exact reduction.
Let $w_u=|A_u|^2-J_u$, let $\mathcal N$ be the finite plain annulus,
and define
\[
 T_{u,v}(N)=\sum_{n\in\mathcal N}\psi_u(n)\overline{\psi_v(n)},
 \qquad B_2=\sum_{n\in\mathcal N}|b_n|^2\ll N.
\]
For the doubly off-diagonal sum,
\begin{align}
 |\mathcal E_{\mathrm{both\ off}}|^2
 &\le \frac{B_2^2}{D^2N^2}
   \sum_{n_1\ne n_2}
    \left|\sum_u w_u\psi_u(n_1)\overline{\psi_u(n_2)}\right|^2\notag\\
 &\le \frac{B_2^2}{D^2N^2}
   \sum_{u,v} w_u w_v |T_{u,v}(N)|^2.\label{eq:cauchy-row-pair}
\end{align}
The last equality follows by expanding the square after adding the
nonnegative $n_1=n_2$ terms. Its right hand side is nonnegative even
though some $w_u$ are negative. Bounding its individual terms further
requires absolute values of these real weights.

The new varying character is exactly
$\psi_u(n)\overline{\psi_v(n)}$.
Its local moving exponent is the difference between the valuations
of $u$ and $v$ modulo six, with the union zero mask retained.
After separating the fixed reciprocity phases, its moving modulus is
supported on $\rad((u)(v))$, with norm at most a fixed-height constant
times $U^2$. For generic coprime bulk rows this norm is comparable to
$U^2$. Thus Cauchy changes the modulus legitimately. It does not apply
a local estimate for one four-ideal modulus to a different unfrozen one.
The additive shifts in a correlation estimate for $T_{u,v}(N)$ are
shifts of $n$, reduced at every prime of this row-pair modulus.
Even if the original four ideals were squarefree and pairwise coprime,
the new rows $u,v$ need not be. Their local exponents can be quadratic,
so the doubled-root exceptions found in the local kernel analysis remain
necessary.

An optimistic square-root bound $|T_{u,v}(N)|\ll N^{1/2}U^\rho$
on all nonexceptional row pairs would still not finish this Cauchy route.
Using $\sum_u|w_u|\ll DU$ gives an off-diagonal upper bound
of order $U N^{1/2}$ for the energy. Even granting an additional ideal
inverse fourth-moment bound $\sum_u|A_u|^4\ll D^2U$ to control
exceptional equal-row pairs, and optimistically assuming there are no
other exceptional row pairs, this leads to
\[
 U^{1+m/2+\rho}+U^{m+1/2+\rho}.
\]
At the crossing the first exponent is $1.2041950721$, exceeding
$1.1028180869$. The exponent comparison would require
$r>1-2m/5$. This condition fails at the crossing.
The square-root row-correlation hypothesis and the inverse fourth moment
are not supplied by the existing proof.

\begin{lemma}[Barrier for direct additive differencing]\label{lem:differencing}
Let $G$ be a finite additive group of order $Q$, let
$\theta:G\to\mathbb C$, and let $\alpha:G\to[0,1]$ satisfy
$\sum_z\alpha(z)\le M$.
For a nonempty shift set $\mathcal H$ of cardinality $H$, set
\[
 \Delta=H^{-1}\sum_{h\in\mathcal H}\sum_z
                  |\alpha(z-h)-\alpha(z)|.
\]
Assume $|\theta|\le1$ and assume the complete fourth-moment estimate
with a fixed constant $C\ge1$, increasing it if needed,
\[
 \sum_z\left|\sum_{h\in\mathcal H}\theta(z+h)\right|^4
 \le C(QH^2+Q^{1/2}H^4).
\]
Then
\[
 \left|\sum_z\alpha(z)\theta(z)\right|
 \le \Delta+C^{1/4}M^{3/4}
          (Q^{1/4}H^{-1/2}+Q^{1/8}).
\]
Consequently this derived bound cannot provide a strict improvement
over the trivial bound $M$ by exponent optimization if $Q\ge M^2$.
This is a statement about the bound obtained by this proof, not a
lower bound on the actual character sum.
\end{lemma}
\begin{proof}
Translate the weights to compare the sum with
$\sum_z\alpha(z)H^{-1}\sum_h\theta(z+h)$.
The absolute difference is at most $\Delta$.
Weighted H\"older bounds the fourth power of the latter by
\[
 M^3 H^{-4}\sum_z\alpha(z)
                \left|\sum_h\theta(z+h)\right|^4.
\]
Discard $\alpha(z)\le1$, apply the assumed complete fourth moment,
and take a fourth root. The second term in this upper bound is at least $M$
when $Q\ge M^2$.
\end{proof}

The complete fourth-moment hypothesis is deliberately conditional.
It would follow from suitable collision counting and a complete
four-distinct-shift Weil bound, followed by absolute summation over the
shift quartets. The local generic candidate is
$|C_P+1|\le2\sqrt{\N P}$, not $|C_P|\le2\sqrt{\N P}$.
CRT collision primes require the exceptional ideal factor retained in
the independent local-kernel note. No uniform collision count for our
moving shift sets is proved here.

Even granting that favorable complete fourth moment and setting
$\Delta=0$, applying Lemma~\ref{lem:differencing} at the original
four-ideal modulus gives $M\asymp U$, $Q\asymp U^{2s}$.
Its unavoidable bound term has exponent $3/4+s/4>1$ when $s>1$.
Applying it instead to \eqref{eq:cauchy-row-pair} gives
$M\asymp U^m$, $Q\asymp U^2$ and exponent $3m/4+1/4>m$ for $m<1$.
The analogous long-factor Cauchy route has exponent
$3r/4+1/4>r$ when $r<1$.
Thus direct additive Cauchy and four-shift absolute completion supplies
no gain at either actual crossing length. The sixth-power-free physical
row indicator additionally has no proved small translation error.
Restricting the original quadruple to squarefree pairwise coprime ideals
does not cover every excluded sector. Shared-prime, nonresonant products
can have both $d_1\ne d_2$ and $n_1\ne n_2$.
Their inactive local primes still impose zero masks.
The controlled strips and resonant bound alone do not dispose of them.
Improving this route requires preserving more coefficient structure,
gaining cancellation between completed shift quartets, or a different
amplification. A local four-shift theorem alone is insufficient.

\paragraph{The remaining falsifiable target.}
Remove from \eqref{eq:typeII} the strips $d_1=d_2$ and $n_1=n_2$,
and remove pairs $d_1n_1,d_2n_2$ that are sixth-equivalent outside
the specified fixed prime set. For fixed common shifts and each total
common mark at most two, prove a bound for the absolute value of the
remaining signed sum
\begin{equation}\label{eq:open-target}
 |\mathcal E_{\mathrm{both\ off,nonres}}|
 \ll (1+H)^J U^{1+5(s-1)/6-\eta+\rho}.
\end{equation}
The region must satisfy $s\ge1+\kappa$ with
$0<\eta<5\kappa/6$, so that the proved $U^{1+\rho}$ pieces fit.
The constants must use fixed smooth seminorms, bounded length ranges,
and the source height controls. The coefficient signs $\mu(d_1)\mu(d_2)$,
the common zero masks, and every fixed reciprocity phase remain in the sum.
No arbitrary independently shifted coefficient sequence is permitted.
Sobolev then supplies the requested moving-shift joint moment.
Equation~\eqref{eq:open-target} is not proved here.

At the original numerical crossing
$\delta=0.3867$, $x=1/2$, the source balance gives
\[
 t=1.1233817043,\quad r=0.7149915600,\quad m=0.4083901443,
 \quad 1+5(t-1)/6=1.1028180869.
\]
Thus a new theorem must beat $U^{1.1028180869}$ there, with a quantitative
positive saving. The complete-residue and controlled-strip results have
exponent one and do fit this budget. They leave a concrete doubly
off-diagonal arithmetic sum rather than a generic coefficient assertion.
At that crossing, a uniform short-count saving $g$ is discounted by about
$0.75324$ after retuning the detector cutoff against the untouched long
branch. Coverage only at this single crossing is insufficient for a new
global boundary, since the active crossing moves as the geometry changes.

\paragraph{Reproducible count.}
Run \texttt{uv run research/joint\_moments.py}.
The script counts actual Eisenstein ideals by norm, using split,
inert, and ramified rational primes. It checks the exact decomposition
$K=A B^6$ at every norm through $524287$.
It then counts ordered sixth-equivalent pairs in $[X,2X)$ exactly.
At $X=65536$ it finds $39631$ resonant pairs and $39629$ diagonal pairs.
At $X=262144$ it finds $158502$ resonant pairs and $158498$ diagonal pairs.
The ratios are about $0.6046X$. These finite counts support the proved
linear counting scale. They do not test the physical incomplete kernel.

\paragraph{Source references.}
All paths below are relative to the upstream \texttt{lean/} directory.

The actual dyadic definition is in
\path{OAI/NumberTheory/DirichletL/Hecke/Dyadic.lean}.

The cutoff and weighted fiber are in
\path{Hecke/DetectorWitnessArithmetic.lean}, lines 11--73.

The exact complete ideal divisor sum is
\path{Arithmetic/IdealMobius.lean}, lines 538--575.

The row character specification is
\path{Hecke/InverseAmplificationFamily.lean}, lines 12--31.

The ideal-row coefficient equality is
\path{Hecke/RowClosure.lean}, theorem \path{idealCoeff_eq_row}.

The local row symbol and product are in
\path{Eisenstein/CompactEnergyFamilies.lean}, lines 466--487.

Canonical sextic and its finite-field Fourier identities are in
\path{GaussSum/FiniteFourier.lean} and
\path{Arithmetic/EisensteinEmbedding.lean}, lines 420--478.

The existing inverse rowwise energy is
\path{Hecke/InverseAmplificationRowwise.lean},
theorem \path{no_slot_rowwise_endpoint}.

The squarefree divisor bound is
\path{Arithmetic/SquarefreeDivisors.lean}, line 130.

The source paper's inverse amplification is at lines 12362--12453
and its plain fourth moment starts at line 12477.

An independent review confirmed the strip, mark, resonant-fiber, and
complete-residue reductions. The four local Lean theorem axiom reports
contain only \texttt{propext}, \texttt{Classical.choice}, and
\texttt{Quot.sound}.
The local kernel agent independently checks the actual finite-field
Fourier phases and provides formal zero-mask lemmas. Those checks do not
prove \eqref{eq:open-target}.

\end{document}
